\documentclass[10pt,a4paper]{article}

% Configurações de idioma e fonte
\usepackage[utf8]{inputenc}
\usepackage[T1]{fontenc}
\usepackage[brazilian]{babel}
\usepackage[margin=0.75in]{geometry}
\usepackage{hyperref}
\usepackage{enumitem}

% Desabilita a numeração de páginas
\pagestyle{empty}

% Formatação personalizada de seções
\usepackage{titlesec}
\titleformat{\section}{\large\bfseries\uppercase}{}{0pt}{}[\titlerule]
\titlespacing{\section}{0pt}{12pt}{8pt}

\begin{document}

% Cabeçalho
\begin{center}
    {\LARGE \textbf{Kevin Liu Rodrigues}}  \\
    \vspace{2pt}
    Brasília, DF, Brazil $|$ +55 19 99732 6791  $|$ \href{mailto:kkevin.liu@gmail.com}{kkevin.liu@gmail.com} $|$
	\href{http://hdl.handle.net/1843/39095}{http://hdl.handle.net/1843/39095}
\end{center}

% Resumo
\noindent\rule{\textwidth}{0.4pt} Sou um cientista com doutorado em física estatística e computacional com sólida base
quantitativa. Comecei a programar durante o mestrado para fazer a interface com diversos equipamentos
eletrônicos e logo migrei para Python, que utilizo desde então.

Inspirado pelos avanços da inteligência artificial com o AlphaGo, no meu doutorado mudei para a área de física
estatística e computacional, implementando várias simulações de sistemas dinâmicos em C. Isso me levou a uma colaboração
com a Universidade de Michigan, onde pude fornecer uma descrição matemática para alguns sistemas de osciladores que
corroboraram os resultados das simulações.

Essa trajetória me proporcionou uma experiência diversificada em ciência de dados, machine learning e desenvolvimento de
software: em 2018, entrei na Olivia AI para desenvolver análise de sentimento e um chatbot antes do advento dos LLMs. Em
2020, fui contratado pela Rappi para desenvolver modelos de ``lifetime value" e métricas de usuários para auxiliar na
expansão da empresa na América do Sul e, finalmente, em 2021, entrei na Spacetime Labs (agora Moray AI) para trabalhar
com machine learning aplicado à agricultura de precisão e manejo de safras por meio de dados climáticos e de sensores.

% Experiência
\section{Experience}

\textbf{Moray} (2021 -- Present)

\textit{Data Scientist and Engineer} 
\begin{itemize}[noitemsep, topsep=0pt]
    \item Machine learning and predictive models for precision agriculture and crop management.
    \item Data pipelines for gathering and processing spatio-temporal data from sources such as NASA, Copernicus, ECMWF, etc.
	\item Time series analysis and forecast.
    \item Machine learning pipelines from training to deployment using mainly MLflow and AWS.
\end{itemize}

\vspace{6pt}

\textbf{Rappi} (2020 -- 2021)

\textit{Data Scientist} 
\begin{itemize}[noitemsep, topsep=0pt]
	\item Development of user life-time value models and user metrics.
	\item Large dataset processing and visualizations including spatio-temporal maps and dashboards.
\end{itemize}

\vspace{6pt}

\textbf{Olivia AI} (2018 -- 2018)

\textit{Data Scientist} 
\begin{itemize}[noitemsep, topsep=0pt]
	\item Development of a chatbot in the very early stages of AI, experimenting with sentiment analysis, string distance metrics and word embeddings.
	\item Classification models for text, including Neural Networks, SVM's and ensemble methods.
\end{itemize}

\vspace{6pt}

\textbf{Federal University of Minas Gerais} (2015 -- 2016)

\textit{Researcher and Engineer} 
\begin{itemize}[noitemsep, topsep=0pt]
	\item Development of a high power near-field microscope for Hewlett-Packard (HP) under coordination of professor Gilberto Medeiros
\end{itemize}

\vspace{6pt}

\textbf{University of Campinas} (2010 -- 2012)

\textit{Scientific initiation} 
\begin{itemize}[noitemsep, topsep=0pt]
	\item Development of pico Faraday detection electronics for capacitive sensing.
	\item Development of electrostatic lensing for ion beam sampling and focusing.
\end{itemize}

\vspace{6pt}

\textbf{TRIU} (2012 -- 2014)

\textit{Volunteer math teacher} 
\begin{itemize}[noitemsep, topsep=0pt]
	\item Math classes for low income students aiming to join top universities.
\end{itemize}

% Educação
\section{Education}

\textbf{Doctorate in statistical physics | UFMG} and \textbf{UMICH} (2016 -- 2021)

(Advisors: Ronald Dickman and Kevin Wood)

\vspace{4pt}

\textbf{Master of science in applied physics | University of Campinas} (2012 -- 2015)

\vspace{2pt}

\textbf{Bachelor of science in physics | University of Campinas} (2008 -- 2012)

% Competências
\section{Skills}
\begin{itemize}[label={}, leftmargin=0pt, noitemsep]
	\item \textbf{Programming languages:} Python, SQL, C, C++, TCL/TK, JavaScript
	\item \textbf{Tools:} MLflow, Docker, Amazon Web Services (AWS), Airflow, Kubernettes, Parametric Modelling (CAD)
    \item \textbf{Languages:} Portuguese (native), English (fluent spoken/written) 
\end{itemize}

\end{document}
